\chapter{From holomorphic brackets to bulk fields}
\label{rc:poisson-sigma}
\label{pcsm:chapter}
\label{pcsm:section}

A source changes the motion of a field through its Poisson bracket.
For a current, this response can depend on derivatives of the source.
The coefficients of those derivatives have a second interpretation:
they are couplings in a field theory with one additional real direction.
The Poisson vertex sigma model of Khan and Zeng~\cite{pcsm:KZ}
makes this construction for a free algebra of differential polynomials.
The Jacobi identity becomes the classical consistency equation for the
bulk interaction.

We shall follow one coefficient through the construction.
A current whose response is \(\kappa\partial\epsilon\), for a source
\(\epsilon\), gives a quadratic bulk action.
Integrating its normal fields produces the Heisenberg boundary algebra
and a supported action on boundary observables.
The action contains a quantum contraction that ordinary restriction
would lose.
Field-dependent response coefficients then lead to affine gauge theory,
Virasoro fields and line transport.

\section{A current and the response to its source}
\label{pcsm:sec:source-response}

Consider an elastic field on the circle
\(s\in\mathbb R/(2\pi\mathbb Z)\).
Its displacement \(q(s)\) and momentum density \(p(s)\) are smooth
and periodic. Set the mass density and tension to one and take the energy
\[
 E=\frac12\int_0^{2\pi}\bigl(p^2+(\partial_s q)^2\bigr)\,ds.
\]
There is no on-site restoring force.

For functionals with variational derivatives, the canonical bracket is
\[
 \{F,G\}_0=\int_0^{2\pi}
 \left(\frac{\delta F}{\delta q}\frac{\delta G}{\delta p}
       -\frac{\delta F}{\delta p}\frac{\delta G}{\delta q}\right)ds.
\]
All integrations by parts use periodicity.
If\phantomsection\label{pcsm:def:physical-time} \(\tau\) denotes physical time, Hamilton's equations
\(\partial_\tau g=\{g,E\}_0\) give
\(\partial_\tau q=p\) and \(\partial_\tau p=\partial_s^2q\).
Consequently momentum and strain have the traveling combination
\begin{equation}\label{pcsm:eq:elastic-current}
 J=\frac{p-\partial_s q}{\sqrt2},\qquad
 \partial_\tau J=-\partial_sJ.
\end{equation}

Couple a smooth periodic source \(\epsilon\) by replacing \(E\) with
\(E-F_\epsilon\), where
\[
 F_\epsilon=\int_0^{2\pi}\epsilon(s)J(s)\,ds.
\]
The change in the rate of an observable \(g\) is then
\(\{F_\epsilon,g\}_0\).
We call this Hamiltonian operation the source response.
The same convention applies at each time to a time-dependent source.

Varying \(F_\epsilon\) gives
\[
 \{F_\epsilon,q(t)\}_0=-\frac{\epsilon(t)}{\sqrt2},\qquad
 \{F_\epsilon,p(t)\}_0=\frac{\epsilon'(t)}{\sqrt2},
\]
where the prime differentiates the spatial argument.
The two contributions to the current therefore add:
\begin{equation}\label{pcsm:eq:current-source-response}
 \{F_\epsilon,J(t)\}_0=\epsilon'(t).
\end{equation}
A constant source has zero response on \(J\), but changes the rate of
the zero mode of \(q\). The current alone does not detect the whole
Hamiltonian action.

The\phantomsection\label{pcsm:def:periodic-contact} unsmeared bracket fixes the contact sign.
Let \(\delta(s-t)\) be the periodic Dirac distribution, normalized by
\(\int_0^{2\pi}\epsilon(s)\delta(s-t)\,ds=\epsilon(t)\).
Since \(\{q(s),p(t)\}_0=\delta(s-t)\),
\begin{align*}
 \{J(s),J(t)\}_0
 &=-\frac12\left(\{p(s),\partial_tq(t)\}_0
                  +\{\partial_sq(s),p(t)\}_0\right)\\
 &=\frac12(\partial_t-\partial_s)\delta(s-t)
   =\partial_t\delta(s-t).
\end{align*}
Smearing the first position recovers
\eqref{pcsm:eq:current-source-response}.
For \(\kappa>0\), the rescaled real current
\(J_\kappa=\sqrt\kappa J\) has response \(\kappa\epsilon'\).

\subsection{Jets and local brackets}

A polynomial measurement at a point uses finitely many derivatives
of the field there. Write
\[
 u_n[J_\kappa]=\left.\partial_s^nJ_\kappa(s)\right|_{s=0},
 \qquad n\geq0.
\]
These jet coordinates generate \(\mathbb C[u_0,u_1,\ldots]\).
Moving the measurement point induces the total derivative
\(\partial u_n=u_{n+1}\).

Equivalently, a formal profile has the Taylor series
\(u(\zeta)=\sum_{n\geq0}u_n\zeta^n/n!\).
Here \(\zeta\) is position, whereas each \(u_n\) is a coordinate
on profiles. No convergence is required for the polynomial algebra.

Encode\phantomsection\label{pcsm:def:contact-translation} differentiation of a contact in its second position by an
even formal variable \(\lambda\).
Differentiating its first position gives \(-\lambda\).
Differentiating its second position also differentiates the coefficient
of the contact. Thus a local bracket must obey
\[
 \{\partial f_\lambda g\}=-\lambda\{f_\lambda g\},\qquad
 \{f_\lambda\partial g\}=(\partial+\lambda)\{f_\lambda g\}.
\]

Products of measurements require the two Leibniz identities
\begin{align*}
 \{f_\lambda gh\}&=\{f_\lambda g\}h+\{f_\lambda h\}g,\\
 \{fg_\lambda h\}&=\{f_{\lambda+\partial}h\}_{\to}g
                    +\{g_{\lambda+\partial}h\}_{\to}f.
\end{align*}
The arrow means that the derivatives introduced by
\(\lambda+\partial\) act on the following factor.
All jet coordinates in this paragraph are even.

The current has \(\{u_\lambda u\}=\kappa\lambda\), with \(u=u_0\).
Translation determines its higher jets:
\begin{equation}\label{pcsm:eq:current-jet-response}
 \{u_m{}_{\lambda}u_n\}=(-1)^m\kappa\lambda^{m+n+1}.
\end{equation}
Leibniz then determines every polynomial bracket.
The constant coefficient makes the generator Jacobi identity zero,
and the opposite contact derivatives give skewsymmetry.
The master formula below extends these identities to all polynomials.

The same formulas define a family
\[
 A_{\mathrm{cur}}=\mathbb C[\kappa][u_0,u_1,\ldots],
 \qquad \partial\kappa=0.
\]
Its zero-level fiber retains every polynomial jet and has zero bracket.
Rescaling the mechanical current described its positive real fibers.
Specialization of this algebra defines the family for every complex
\(\kappa\).

\subsection{The Hamiltonian differential operator}

For\phantomsection\label{pcsm:def:holomorphic-coordinate} a solution admitting continuation to \(\tau=-i\tau_E\), put
\(z=s+i\tau_E\).
Equation~\eqref{pcsm:eq:elastic-current} becomes
\(\partial_{\bar z}J=0\).
We may also complexify the polynomial bracket directly, without
continuing each smooth profile. The total derivative then represents
holomorphic translation.

For\phantomsection\label{pcsm:def:field-jet-algebra} several fields, take
\[
 A=\mathbb C[u^a_n:1\leq a\leq r,\ n\geq0],\qquad
 \partial u^a_n=u^a_{n+1},\qquad u^a=u^a_0.
\]
A \emph{Poisson vertex algebra} structure on \(A\) is a bilinear map
\(\{-_\lambda-\}:A\otimes A\to A[\lambda]\)
satisfying the two translation and two Leibniz identities above,
skewsymmetry and Jacobi:
\begin{align*}
 \{f_\lambda g\}&=-\{g_{-\lambda-\partial}f\},\\
 \{f_\lambda\{g_\mu h\}\}-\{g_\mu\{f_\lambda h\}\}
 &=\{\{f_\lambda g\}_{\lambda+\mu}h\}.
\end{align*}
If \(\{g_\mu f\}=\sum_n b_n\mu^n\), the substitution in the
first identity is \(\sum_n(-\lambda-\partial)^n b_n\).
Polynomial dependence on \(\lambda\) expresses finite differential
order for each pair of fields.

Write\phantomsection\label{pcsm:def:hamiltonian-response} the generator brackets as
\[
 \{u^a_\lambda u^b\}=H^{ba}(\lambda),\qquad
 H^{ab}(\partial)=\sum_{n=0}^N
                 H_n^{ab}(u,\partial u,\ldots)\partial^n.
\]
The transpose records which field responds: a source \(\epsilon_b\)
coupled to \(u^b\) induces
\(\delta_\epsilon u^a=H^{ab}(\partial)\epsilon_b\).
Thus \(H\) maps a column of source covectors to field variations.
Repeated indices below are summed from \(1\) to \(r\).

Translation and Leibniz give the explicit extension
\[
 \{f_\lambda g\}
 =\sum_{a,b,m,n}
 \frac{\partial g}{\partial u^b_n}(\lambda+\partial)^n
 H^{ba}(\lambda+\partial)(-\lambda-\partial)^m
 \frac{\partial f}{\partial u^a_m}.
\]
Every operator acts on what follows it. Each sum is finite because
\(f,g\) use finitely many jets and \(H\) has finite order.

The formal adjoint reverses composition and sends
\(\partial\) to \(-\partial\).
The extended bracket is skew precisely when
\(H^{ab}=-(H^{ba})^*\). Under this condition its Jacobi identity
is equivalent to Jacobi on generator triples
\cite{pcsm:BDK}*{Theorem 1.15 and Proposition 1.16}.
Such an \(H\) is a Hamiltonian differential operator.

For the current, \(H=\kappa\partial\).
Allowing its coefficients to depend on fields introduces interactions.
The bulk construction uses this same operator in its action, so the
identity controlling successive source responses also controls the
gauge interaction.

\section{Fields and the classical action}

Work\phantomsection\label{pcsm:def:partial-forms} on \(M=\mathbb R_{\geq0,t}\times\mathbb C_z\),
with \(z=x+iy\) and orientation \(dt\wedge dx\wedge dy\).
Partial forms are smooth complex forms generated by \(dt,d\bar z\).
Their complex is \(\Omega_{\mathrm{HT}}^\bullet(M)\), with
\[
 D=dt\,\partial_t+d\bar z\,\partial_{\bar z},\qquad
 \partial=\partial_z.
\]
The holomorphic derivative has degree zero and commutes with \(D\).
Fields and variations are smooth up to the boundary.
Their compact supports may meet \(t=0\).

For compactly supported partial forms, fix the traces
\[
 \tau_M(a)=\frac1{2\pi i}\int_M dz\wedge a,
 \qquad
 \tau_{\mathbb C}(b)=\frac1{2\pi i}\int_{\mathbb C}dz\wedge b.
\]
They select partial form degree two and antiholomorphic degree one,
respectively. The second uses the complex orientation.
With\phantomsection\label{pcsm:def:boundary-inclusion} \(i:\mathbb C\hookrightarrow M\) the boundary inclusion,
Stokes' formula in these conventions is
\(\tau_M(Da)=\tau_{\mathbb C}(i^*a)\).
The sign includes the outward boundary orientation and the position
of \(dz\) in the trace.

\subsection{The scalar action fixes the gauge sign}

Let \(\phi\) be a scalar and
\(a=a_tdt+a_{\bar z}d\bar z\) a partial one-form, both compactly
supported. For \(\kappa\in\mathbb C\), set
\[
 S^{\mathrm{cl}}_\kappa(\phi,a)
 =\tau_M\left(aD\phi+\frac\kappa2a\partial a\right).
\]
The first term couples the current coordinate to partial transport.
The second inserts the measured derivative response.
Its coefficient of \(dt\wedge d\bar z\) is
\[
 a_t\partial_{\bar z}\phi-a_{\bar z}\partial_t\phi
 +\frac\kappa2(a_t\partial a_{\bar z}-a_{\bar z}\partial a_t).
\]

Varying and integrating by parts gives
\[
 \delta S^{\mathrm{cl}}_\kappa
 =\tau_M\bigl(\delta a(D\phi+\kappa\partial a)
                         +(Da)\delta\phi\bigr)
       -\tau_{\mathbb C}(i^*a\,\delta\phi).
\]
Thus the equations are \(D\phi+\kappa\partial a=0\) and \(Da=0\).
The condition \(i^*a=0\) cancels the boundary variation and leaves
\(\phi\) free there.

For\phantomsection\label{pcsm:def:scalar-gauge} a compactly supported parameter \(\epsilon\) vanishing on the
boundary, define
\[
 \delta_\epsilon\phi=\kappa\partial\epsilon,
 \qquad \delta_\epsilon a=-D\epsilon.
\]
Substitution in the variation cancels the two terms containing
\(\partial Da\), after integration by parts. The remaining term
vanishes by \(D^2=0\). The same identities preserve both equations.
The coefficient of the source response has therefore fixed a gauge
symmetry as well as the action.

\subsection{The full BV fields}

The Batalin--Vilkovisky description puts gauge transformations,
equations of motion and their identities into one differential.
For the Hamiltonian operator \(H\), collect the scalar fields,
degree-one ghosts and their antifields into superfields
\[
 \Phi^a=\phi^a+dt\,\phi^a_t+d\bar z\,\phi^a_{\bar z}
                  +dt\,d\bar z\,\phi^a_{t\bar z},
\]
\[
 \eta_a=c_a+dt\,\eta_{a,t}+d\bar z\,\eta_{a,\bar z}
                  +dt\,d\bar z\,\eta_{a,t\bar z}.
\]

Their total degrees are zero and one.
The coordinate degrees in the first row are \(0,-1,-1,-2\).
Those in the second are \(1,0,0,-1\).
The physical fields are \(\phi^a,\eta_{a,t},\eta_{a,\bar z}\),
and \(c_a\) is the gauge ghost. The remaining coordinates are
antifields, paired across complementary partial form degrees.

The two-form
\(\omega_{\mathrm{BV}}=\tau_M(\delta\Phi^a\delta\eta_a)\)
has degree \(-1\). Its inverse gives the degree-one BV bracket on
local functionals. Replace the scalar operator in the action by \(H\):
\begin{equation}\label{pcsm:action}
 S_H=\tau_M\left(\eta_aD\Phi^a
          +\frac12\eta_aH^{ab}(\partial;\Phi)\eta_b\right).
\end{equation}
Both densities have total degree two, so \(S_H\) has degree zero.
For \(H=\kappa\partial\), its physical component is
\(S^{\mathrm{cl}}_\kappa\).

\begin{proposition}[Classical model]\label{pcsm:classical}
For a Hamiltonian differential operator \(H\), the action
\eqref{pcsm:action} obeys the classical master equation
\(\{S_H,S_H\}_{\mathrm{BV}}=0\) on fields with boundary condition
\begin{equation}\label{pcsm:boundary}
 i^*\eta=0.
\end{equation}
This condition on the whole superfield is Lagrangian and is preserved
by the BV differential \(Q=\{S_H,-\}_{\mathrm{BV}}\).
\end{proposition}

\begin{proof}
The kinetic term squares to zero by \(D^2=0\).
Its bracket with the interaction is a total partial derivative.
Its boundary trace contains \(i^*\eta\), so it vanishes.

For the interaction bracket with itself, pair the coefficients with
three independent odd test fields and integrate the holomorphic
derivatives by parts. The coefficient is the generator Jacobi identity
\[
 \{u^a_\lambda\{u^b_\mu u^c\}\}
 -\{u^b_\mu\{u^a_\lambda u^c\}\}
 -\{\{u^a_\lambda u^b\}_{\lambda+\mu}u^c\}=0.
\]
The Hamiltonian-operator identity of
\cite{pcsm:BDK}*{Proposition 1.16} is exactly this calculation after
integration by parts. Hence the master equation holds, and graded
Jacobi gives
\(Q^2=\tfrac12\{\{S_H,S_H\}_{\mathrm{BV}},-\}_{\mathrm{BV}}=0\).

Variation of the action has boundary term
\(-\tau_{\mathbb C}(\eta_a\delta\Phi^a)\).
The condition \eqref{pcsm:boundary} sets one of the complementary
summands of the nondegenerate boundary pairing to zero and leaves
the other unrestricted. It is Lagrangian.
The interaction in \(Q\eta\) is quadratic in \(\eta\) and its
holomorphic derivatives. Its tangential pullback vanishes when
\(i^*\eta=0\), and the kinetic boundary differential
\(\bar\partial\) preserves the same condition.
\end{proof}

This constructs a local classical theory of finite holomorphic
differential order \cite{pcsm:KZ}*{\S\S2.2 and 4.1}.
Quantizing a field-dependent \(H\) requires a further construction.

Already among constant operators there are different boundary fields.
A constant order-zero skew matrix of rank \(2s\) has \(s\)
two-by-two skew blocks and a zero block of dimension \(r-2s\).
Each nonzero block gives a beta--gamma pair
\(\{\beta^i_\lambda\gamma_j\}=\delta^i_j\), while the zero block
gives commutative differential currents. Only a nondegenerate matrix
gives a pure beta--gamma system. The positive-order scalar operator
\(\kappa\partial\) will instead give a double pole.

\section{The cotangent coisson complex}
\label{pcsm:cotangent}

The same Hamiltonian also acts on a purely holomorphic polynomial
complex. Adjoin odd variables \(c_{a,n}=\partial^nc_a\) of degree
one to the degree-zero algebra \(A\), and put
\[
 \mathcal P_H=A\otimes\Lambda(c_{a,n}:1\leq a\leq r,\ n\geq0).
\]

Its Schouten \(\lambda\)-bracket has degree \(-1\), with generator
values
\[
 \{c_a{}_\lambda u^b\}_{\mathrm{Sch}}=\delta_a^b,
 \qquad \{u^b{}_\lambda c_a\}_{\mathrm{Sch}}=-\delta_a^b,
\]
and zero on the other pairs. Extend by sesquilinearity and graded
Leibniz. The integrated Hamiltonian density
\(\tfrac12c_aH^{ab}(\partial)c_b\) determines the differential
\begin{align}\label{pcsm:cotangent-differential}
 Q_Hu^a&=H^{ab}(\partial)c_b,\notag\\
 Q_Hc_a&=\frac12\sum_{n\geq0}(-\partial)^n
   \left(c_b\frac{\partial H^{bc}(\partial)}{\partial u^a_n}c_c\right),
 \qquad [Q_H,\partial]=0.
\end{align}
The differentiated operator still acts on \(c_c\).
With the bracket convention just chosen, \(Q_H\) is minus the
Hamiltonian action of that density. Its square vanishes by the
same variational Jacobi identity.

Thus \(H\) has supplied two complexes independently: the BV fields
on the half-space and \(\mathcal P_H\) on holomorphic jets.
Their brackets have different degrees because the BV trace integrates
two partial form directions. Identifying their observables requires
a map. For the Gaussian and affine models we shall use an explicit
contraction of the full component jets.
Quotienting \(\mathcal P_H\) by total holomorphic derivatives is
another operation. It will change the Virasoro cohomology calculation.

\section{The Gaussian theory on open sets}

For \(H=\kappa\partial\), the complete BV action is quadratic:
\begin{equation}\label{pcsm:freeaction}
 S_\kappa=\tau_M\left(\eta D\Phi+\frac\kappa2\eta\partial\eta\right).
\end{equation}
Keep \(\kappa\) as a parameter in
\(\mathbb k_\kappa=\mathbb C[\kappa]\).
The regular observables and collar maps will be defined over this ring,
so they specialize at every complex level.

A\phantomsection\label{pcsm:def:parameter-field-complex} tangent presentation of the fields is
\[
 E=\bigl(\Omega_{\mathrm{HT}}^\bullet e_\Phi
          \oplus\Omega_{\mathrm{HT}}^\bullet e_\eta\bigr)
             \otimes\mathbb k_\kappa,
 \qquad |e_\Phi|=0,\quad |e_\eta|=-1,
\]
with differential
\[
 Q(\alpha e_\Phi)=D\alpha e_\Phi,\qquad
 Q(\alpha e_\eta)=D\alpha e_\eta
             +(-1)^{p+1}\kappa\partial\alpha e_\Phi,
 \quad p=\deg\alpha.
\]
These are tangent degrees. The coordinate degrees of the superfields
are dual to them. The sign makes the two off-diagonal terms in
\(Q^2\) cancel. The \(e_\Phi\) summand is a subcomplex, and the
associated graded symbol consists of two elliptic mixed
de Rham--Dolbeault complexes. Impose \(i^*\eta=0\) at the boundary.

The degree-minus-one pairing is fixed on components by
\[
 B_M(\alpha e_\Phi,\beta e_\eta)=\tau_M(\alpha\beta),\qquad
 B_M(\alpha e_\eta,\beta e_\Phi)
       =-(-1)^{\deg\beta}\tau_M(\alpha\beta),
\]
and is zero on the two equal summands.
The boundary Green pairing uses the same formulas with
\(\tau_{\mathbb C}\).
This convention recovers the quadratic action
\eqref{pcsm:freeaction}.

For\phantomsection\label{pcsm:def:linear-regular-observables} an open \(U\subset M\), let \(E_c(U)\) denote compactly
supported smooth fields, with the boundary condition whenever
\(U\) meets \(t=0\).
Linear regular observables are modeled by \(C(U)=E_c(U)[1]\).
If \(s\) denotes suspension, so \(|s\xi|=|\xi|-1\), put
\[
 d(s\xi)=-sQ\xi,\qquad
 b(s\xi,s\zeta)=(-1)^{|\xi|+1}B_M(\xi,\zeta).
\]

The form \(b\) has degree one and is graded symmetric.
On a symmetric product, let \(\Delta_b\) sum its unordered pair
contractions, with the Koszul signs. In particular it kills constants
and linear factors, and \(\Delta_b(vw)=b(v,w)\).
Compatibility of \(Q\) with the pairing gives
\(d\Delta_b+\Delta_bd=0\). Odd pair contractions give
\(\Delta_b^2=0\).

Set\phantomsection\label{pcsm:def:quantum-coefficient-ring} \(R=\mathbb k_\kappa[\hbar]\), with \(|\hbar|=0\).
The quantum observable complex is
\begin{equation}\label{pcsm:quantobs}
 \mathcal O^q(U)=
 \left(\bigoplus_{n\geq0}\operatorname{Sym}_b^n C(U)[\hbar],
                         d+\hbar\Delta_b\right).
\end{equation}
At fixed field number \(n\), the subscript \(b\) means the
completed bornological tensor product of compact smooth kernels,
followed by graded symmetrization over \(\mathbb k_\kappa\).
At fixed compact support and parameter degree, the smooth coefficient
spaces are nuclear Fr\'echet and this tensor is the completed
projective tensor.
We then take the bornological inductive limit over supports
\cite{pcsm:GRW}*{Appendix B}.
The sum over field number and the parameter coefficients are algebraic.

Regularity is essential to the contraction.
Evaluation of a field at a point has a delta-distribution kernel.
It is absent from \eqref{pcsm:quantobs}.
Here restriction of a smooth joint kernel to a diagonal is defined,
and integration contracts the two slots.
Extending these operations to local composite distributions requires
renormalized observables.

For\phantomsection\label{pcsm:def:hbar-completion} finite precision write \(R_N=R/(\hbar^{N+1})\).
If \(\mathcal V=\bigoplus_n\operatorname{Sym}_b^n C(U)\), its
\(\hbar\)-adic completion is
\[
 \mathcal V[[\hbar]]=\varprojlim_N
                  \mathcal V\otimes_{\mathbb k_\kappa}R_N.
\]
Each coefficient uses finitely many field numbers.
Their bound may grow with the power of \(\hbar\), as it does in
\(\sum_{m\geq0}\hbar^m x^m\).
Thus this completion generally exceeds
\(\bigoplus_n(\operatorname{Sym}_b^n C(U))[[\hbar]]\).
The finite-precision constructions always use the algebraic sum.

\subsection{Factorization and the normal direction}

Disjoint supports permit multiplication before any choice of
propagator. Given pairwise disjoint opens \(U_i\subset V\), extend
the kernels by zero and concatenate their variables. This defines
\[
 \mathcal O^q(U_1)\otimes_R\cdots\otimes_R\mathcal O^q(U_m)
          \longrightarrow\mathcal O^q(V).
\]
There are no cross contractions on disjoint supports, so the map
commutes with \(d+\hbar\Delta_b\).
The maps are continuous at fixed field numbers and compose by
extension by zero. The structure they express has a useful formulation
in any complex boundary dimension.

\begin{definition}\label{fs:def:bulk-boundary}
Fix a symmetric monoidal category of differentiable cochain complexes.
A \emph{factorization algebra} assigns a complex
\(\mathcal F(U)\) to each open set and maps
\(\bigotimes_i\mathcal F(U_i)\to\mathcal F(V)\) for disjoint
\(U_i\subset V\). They respect composition, permutations and units.
Finite disjoint unions give tensor products and the empty open gives
the monoidal unit. The augmented homotopy \v{C}ech diagram must be an
equivalence for each Weiss cover, a cover containing every finite
subset in some member.
For smooth regular observables we use complete bornological complexes
with their completed bornological tensors.

A \emph{bulk--boundary system} on
\(\mathbb R_{\geq0}\times\mathbb C^d\) is such a factorization
algebra \(\mathcal B\).
It has differentiable complex translations and local real translations
in the interior, compatible with its operations. Writing \(d_{\mathcal B}\) for its differential,
holomorphicity is specified by translation-invariant factorization
derivations \(\eta_j\) of degree \(-1\) satisfying
\[
 [d_{\mathcal B},\eta_j]=\partial_{\bar z_j},\qquad
 [\eta_i,\eta_j]=0.
\]
The brackets are graded commutators.
Normal local constancy requires the maps
\(\mathcal B(I\times U)\to\mathcal B(I'\times U)\)
to be quasi-isomorphisms for nonempty intervals
\(I\subset I'\subset(0,\infty)\), and likewise for inclusions
of boundary collars \([0,\varepsilon)\times U\).
\end{definition}

The\phantomsection\label{pcsm:def:boundary-algebra} boundary algebra is \(\mathcal A(U)=
\mathcal B([0,\varepsilon)\times U)\). Collar inclusions give
coherent quasi-isomorphisms as \(\varepsilon\) changes.
The real dimensions are \(2d\) at the boundary and \(2d+1\) in
the bulk. A translation-invariant system on the whole real line needs an
extension with the same holomorphic and normal structures.
The extension must be locally equivalent to the interior restriction.

For the free elliptic theory, regular observables satisfy this descent
and the normal comparison
\cite{pcsm:GRW}*{Theorems 4.1--4.2}.
We now construct that comparison and compute the induced boundary
pairing. This will fix the current level rather than merely identify
its underlying polynomial states.

\section{Integrating the normal fields}

Fix\phantomsection\label{pcsm:def:collar-profile} an open \(V\subset\mathbb C\), choose \(\varepsilon>0\),
and let \(\rho\in C_c^\infty([0,\varepsilon))\) have integral one.
Put \(F(t)=\int_0^t\rho(s)\,ds\).
The compact tangential fields split into the boundary condition and
its complement:
\[
 E_\partial=L\oplus G,\qquad
 L=\Omega_c^{0,\bullet}(V)e_\Phi\otimes\mathbb k_\kappa,
 \quad G=\Omega_c^{0,\bullet}(V)e_\eta\otimes\mathbb k_\kappa.
\]
In\phantomsection\label{pcsm:def:collar-splitting} this splitting the tangential differential has the form
\[
 d_\partial(\ell,g)=(d_L\ell+Q_{\mathrm{rel}}g,d_Gg),
 \qquad d_LQ_{\mathrm{rel}}+Q_{\mathrm{rel}}d_G=0,
\]
where \(|Q_{\mathrm{rel}}|=1\) and
\(Q_{\mathrm{rel}}(\alpha e_\eta)=
(-1)^{\deg\alpha+1}\kappa\partial\alpha e_\Phi\).

Write a collar field with \(dt\) on the left as
\(\gamma(t)+dt\,\beta(t)\), with \(\gamma_G(0)=0\).
Its differential is
\[
 Q(\gamma+dt\,\beta)=d_\partial\gamma
                +dt(\partial_t\gamma-d_\partial\beta).
\]

Integration of the normal one-form component maps it to \(G[-1]\),
whose differential is \(-d_G\).
A cochain section needs a correction to the bump form:
\begin{equation}\label{pcsm:normalI}
 P(\gamma+dt\,\beta)=\int_0^\infty\beta_G(s)\,ds,
 \qquad Ig=dt\,\rho g+(F-1)Q_{\mathrm{rel}}g.
\end{equation}
Indeed \(Q(dt\,\rho g)\) has the unwanted term
\(-dt\,\rho Q_{\mathrm{rel}}g\). Differentiating \(F-1\)
cancels it. Consequently \(QI=-Id_G\) and \(PI=1\).

The corresponding homotopy is
\[
 h(\gamma+dt\,\beta)=-\int_t^\infty\beta(s)\,ds
       +(1-F(t))\int_0^\infty\beta_G(s)\,ds.
\]
The second term sets its \(G\)-component to zero at the boundary.
The tail integrals vanish above the input and profile supports,
and every map retains compact tangential support.
The fundamental theorem of calculus now gives
\[
 Qh+hQ=1-IP,
 \qquad hI=Ph=h^2=0.
\]
These formulas apply continuously to smooth collar fields, including
their completed kernels. They use no inverse of \(\kappa\).

\subsection{The induced current pairing}

Let\phantomsection\label{pcsm:def:current-identification} \(\mathfrak a(V)=\Omega_c^{0,\bullet}(V)\otimes\mathbb k_\kappa\)
with differential \(\bar\partial\).
Since \(d_G=\bar\partial\) on coefficients, the degree-preserving
identification with \(G[-1]\) requires a sign.
Choose \(\zeta\in\mathbb C\) with \(\zeta^2=-1\) and set
\[
 \jmath_\partial(\alpha)=\zeta(-1)^q\alpha e_\eta
       \quad\text{in }G[-1],\qquad q=\deg\alpha.
\]
Then \((-d_G)\jmath_\partial=\jmath_\partial\bar\partial\).

On shifted linear observables define
\begin{equation}\label{pcsm:eq:current-inclusion}
 i(s\alpha)=sI\jmath_\partial(\alpha),\qquad
 p(s\xi)=s\jmath_\partial^{-1}P\xi,\qquad
 h_C(s\xi)=-s h(\xi).
\end{equation}
The map \(i:\mathfrak a(V)[1]\to E_c(V\times\mathbb R_{\geq0})[1]\)
lands in the chosen collar. The map \(p\) goes in the reverse direction.
They intertwine the unary differentials and satisfy
\[
 pi=1,\qquad dh_C+h_Cd=1-ip,\qquad
 h_Ci=ph_C=h_C^2=0.
\]

The square root \(\zeta\) fixes the pairing as well as the chain
identification. For Dolbeault degrees \(q,r\), only \(q+r=1\)
contributes. The trace convention reads
\[
 \tau_M(dt\,\omega)=-\int_0^\infty\tau_{\mathbb C}(\omega)\,dt.
\]
Substitute the two components of \(I\), integrate the holomorphic
derivative by parts, and use
\(2\int_0^\infty\rho(1-F)\,dt=1\).
The result is
\[
 B_M\bigl(I(\alpha e_\eta),I(\beta e_\eta)\bigr)
       =\kappa\tau_{\mathbb C}(\alpha\partial\beta).
\]

The two \(\jmath_\partial\) factors contribute
\(\zeta^2(-1)^{q+r}=1\), and suspension contributes
\((-1)^{q+1}\).
Define the current cocycle
\begin{equation}\label{pcsm:cocycle}
 \mu_\kappa(\alpha,\beta)
       =-\kappa\tau_{\mathbb C}(\alpha\partial\beta).
\end{equation}
We have proved the normalized identity
\begin{equation}\label{pcsm:eq:current-pairing}
 b\bigl(i(s\alpha),i(s\beta)\bigr)
                    =(-1)^q\mu_\kappa(\alpha,\beta).
\end{equation}

Holomorphic\phantomsection\label{pcsm:def:current-central-extension} integration by parts makes \(\mu_\kappa\) graded
skew, and compatibility with \(\bar\partial\) makes it a cocycle.
Adjoin a central element \(K\) of degree one to the abelian dg Lie
algebra \(\mathfrak a(V)\), with
\([\alpha,\beta]=\mu_\kappa(\alpha,\beta)K\).
In its Chevalley--Eilenberg chains the unary and binary differentials are
\[
 d_1(s\alpha)=-s\bar\partial\alpha,
 \qquad
 d_2(s\alpha\,s\beta)=(-1)^{\deg\alpha}
                           \mu_\kappa(\alpha,\beta)sK.
\]
Specializing\phantomsection\label{pcsm:def:boundary-current-complex} the degree-zero element \(sK\) to \(\hbar\) gives
the boundary current complex
\[
 \mathcal A_\kappa(V)=
 \left(\operatorname{Sym}_b(\mathfrak a(V)[1])[\hbar],
                  d+\hbar\Delta_{\mu^s}\right),
 \qquad \mu^s=i^*b.
\]
Its disjoint-support products are defined by extension by zero.
Equation~\eqref{pcsm:eq:current-pairing} determines the binary
differential, including its suspension sign.

\subsection{The quantum projection}

The linear section preserves the pairing, so its symmetric extension
is already a quantum cochain map.
Ordinary projection loses contractions involving the eliminated fields.
The normal homotopy supplies exactly their correction.

\begin{proposition}[Quantum normal projection]\label{pcsm:quantprojection}
The regular coupled collar observables contract onto
\(\mathcal A_\kappa(V)\) over \(R\).
The maps and homotopy are continuous and polynomial in \(\hbar\)
on each finite field-number component.
They specialize at every level and extend to the stated
\(\hbar\)-adic completion.
\end{proposition}

\begin{proof}
Put\phantomsection\label{pcsm:def:pair-contraction} \(e=ip\) and write \(h=h_C\).
On two shifted factors, average
\(h\otimes1+e\otimes h\) with its graded transpose to obtain
\(H_2\). The tensor differential gives
\(dH_2+H_2d=1-e\otimes e\).
Thus the degree-zero graded-symmetric form \(P_2=-bH_2\) has
Hom differential
\[
 \delta P_2=b-p^*\mu^s.
\]

Let \(\partial_{P_2}\) contract each unordered pair by \(P_2\).
It vanishes on field numbers zero and one, and
\[
 [d,\partial_{P_2}]=\Delta_b-\Delta_{p^*\mu^s}.
\]

Constant-coefficient contraction operators graded-commute.
In particular the even operator \(\partial_{P_2}\) commutes with
both odd Laplacians. For \(T=\exp(\hbar\partial_{P_2})\),
the conjugation series therefore stops after its first commutator:
\[
 T(d+\hbar\Delta_b)T^{-1}=d+\hbar\Delta_{p^*\mu^s}.
\]
It follows that the quantum maps are
\begin{equation}\label{pcsm:quantmaps}
 i^q=\operatorname{Sym}(i),\qquad
 p^q=\operatorname{Sym}(p)\exp(\hbar\partial_{P_2}).
\end{equation}
Since \(H_2\) vanishes on \(i\otimes i\), one has
\(Ti^q=i^q\) and \(p^qi^q=1\).

For the homotopy, split the linear complex as
\(\operatorname{im}i\oplus\ker p\).
Extend \(h\) as a derivation of its symmetric algebra and divide
by the number of \(\ker p\) factors whenever this number is positive.
Set the result \(H_0\) to zero when there are none.
The linear contraction implies
\([d,H_0]=1-\operatorname{Sym}(e)\).
Only retained factors enter \(\Delta_{p^*\mu^s}\), so its graded
commutator with \(H_0\) is zero. Consequently
\[
 H^q=T^{-1}H_0T,\qquad
 (d+\hbar\Delta_b)H^q+H^q(d+\hbar\Delta_b)=1-i^qp^q.
\]

Every contraction lowers field number by two, so the exponentials are
finite on a polynomial. Normal integration, restriction of smooth
joint kernels to diagonals and compact integration are continuous
on each fixed-support kernel space.
Their extensions to the bornological inductive limits give the
asserted continuity. At each power of \(\hbar\) the completed
formulas involve finitely many input coefficients, which proves the
last assertion.
\end{proof}

\section{The boundary current and its stress field}

The full smooth disk complex contains more states than a vertex
algebra's algebraic energy decomposition.
Take the algebraic sum of rotation eigenspaces, called the
\emph{finite-energy sector}.
Its creation modes can be fixed by explicit compact representatives.

Choose a radial annular cutoff \(\chi\), equal to one at the inner
circle and zero at the outer circle. For a holomorphic function on
the annulus, Stokes' formula gives
\[
 \tau_{\mathbb C}((\bar\partial\chi)f)
                  =\operatorname{Res}_0(f,dz).
\]
Hence \(\alpha_n=(\bar\partial\chi)z^{-n}\), for \(n\geq1\), obeys,
for \(m\geq0\),
\[
 \tau_{\mathbb C}(z^m\alpha_n)=\delta_{m,n-1},\qquad
 -\partial\alpha_n=n\alpha_{n+1}
             -\bar\partial(z^{-n}\partial_z\chi).
\]
The last term is compactly supported and exact.

These formulas fix both the mode normalization and translation.

A vertex algebra packages fields with a vacuum and translation.
A state \(v\) has a Laurent field \(Y(v,z)\).
Its application to the vacuum is regular at \(z=0\) with value \(v\).
Translation differentiates fields.
Two fields commute after multiplication by a sufficiently large
power of their separation.
In the present example the double-pole current generates every state.

\begin{theorem}[Heisenberg boundary]\label{pcsm:heis}
The finite-energy vertex algebra of \(\mathcal A_\kappa\) is the
universal Heisenberg vacuum algebra over \(R\), with
\begin{equation}\label{pcsm:ope}
 J(z)J(w)\sim\frac{\hbar\kappa}{(z-w)^2},
 \qquad [J_\lambda J]=\hbar\kappa\lambda.
\end{equation}
This is the exact Gaussian perturbative boundary quantization.
Its vacuum and translation persist at \(\kappa=0\).
\end{theorem}

\begin{proof}
Proposition~\ref{pcsm:quantprojection} identifies the regular collar
complex with the normalized current complex.
For Laurent polynomials \(f,g\), its cocycle is
\[
 \omega_\kappa(f,g)
  =-\kappa\operatorname{Res}_0(f\partial g,dz)
  =\kappa\operatorname{Res}_0((\partial f)g,dz).
\]
Thus \(J_m=\operatorname{Res}_0(z^mJ(z),dz)\) satisfies
\([J_m,J_n]=\hbar\kappa m\delta_{m+n,0}\).
The vacuum representation is explicitly
\[
 V_\kappa=R[x_1,x_2,\ldots],\qquad \mathbf1=1,
 \qquad
 J_{-n}=x_n,\quad
 J_n=\hbar\kappa n\partial_{x_n}\ (n>0),\quad J_0=0.
\]
Each state is a finite polynomial, and \(x_n\) has energy \(n\).

The cutoff representatives above give its translation
\[
 T=\sum_{n\geq1}n x_{n+1}\partial_{x_n},\qquad
 [T,J_m]=-mJ_{m-1}.
\]
The mode commutator is equivalent to the double pole in
\eqref{pcsm:ope}, and implies
\((z-w)^2[J(z),J(w)]=0\).

Normal ordering places positive modes to the right.
Normal-ordered products of derivatives of \(J\) supply the fields
of every polynomial state. For example
\(\partial^aJ(z)\mathbf1|_{z=0}=a!x_{a+1}\).
Their Wick expansions give the vacuum, translation and locality
identities. This is the abelian current extraction of
\cite{pcsm:GRW}*{\S5.3}.

At zero level, all positive modes vanish.
The same carrier becomes the commutative differential vertex algebra
with derivation \(T\). No simple quotient has been taken.

Quadraticity makes Wick contraction exact.
If \(\kappa\eta\partial\eta/2\) is treated as an insertion,
each propagator joins \(\Phi\) to \(\eta\).
Two such insertions cannot join each other.
A boundary pair of \(\Phi\)'s receives exactly one insertion,
whose two attachments cancel its factor \(1/2\).
Thus there is no interaction correction to the current level.
Normalized polynomial correlations do not determine a vacuum
determinant or a framing-dependent partition factor.
\end{proof}

\input{manuscript/constructions/benchmark-boundary-chiral.tex}

\subsection{A quantum term in the composite stress field}

The\phantomsection\label{pcsm:def:composite-stress} current's quadratic composite has a double contraction even
though its level is unchanged.
After inverting \(\kappa\), define
\(L=(2\kappa)^{-1}:J^2:\) over \(R[\kappa^{-1}]\).
Wick's rule counts four single contractions and two double ones:
\[
 :J^2:(z):J^2:(w)\sim
 \frac{4\hbar\kappa}{(z-w)^2}:J(z)J(w):
       +\frac{2\hbar^2\kappa^2}{(z-w)^4}.
\]

Expand the remaining field at \(w\) and use
\(\partial:J^2:=2:(\partial J)J:\). The stress product is
\begin{equation}\label{pcsm:stress}
 L(z)L(w)\sim\frac{\hbar^2/2}{(z-w)^4}
       +\frac{2\hbar L(w)}{(z-w)^2}
       +\frac{\hbar\partial L(w)}{z-w}.
\end{equation}
Equivalently,
\([L_\lambda L]=\hbar(\partial+2\lambda)L+\hbar^2\lambda^3/12\).

The double-contraction graph has two vertices and two edges, hence
one cycle. It changes the composite OPE while leaving \(\kappa\)
fixed. The usual conformal field of central charge one is
\(L/\hbar\), which additionally requires \(\hbar^{-1}\).

\section{Point-local bulk operators}

The supported quantum complex answers how observables in different
regions compose. A classical measurement at one point leads instead
to the polynomial algebra of component jets.
We first identify the gauge theory at nonzero level, then contract
all of its transverse component jets without inverting the level.

Assume\phantomsection\label{pcsm:def:full-exterior-derivative} \(\kappa\ne0\) for the change of variables and write
\(d_M=D+dz\,\partial\) for the full exterior derivative.
The complete abelian superconnection is
\begin{equation}\label{pcsm:csmap}
 A=\eta-\kappa^{-1}\Phi\,dz.
\end{equation}
Exterior multiplication gives
\begin{equation}\label{pcsm:csaction}
 -\frac{\kappa}{2}\frac1{2\pi i}\int_M A\wedge d_MA
       =S_\kappa+\frac12\tau_M D(\Phi\eta).
\end{equation}

For\phantomsection\label{pcsm:eq:connection-three-form} the ordinary three-form component the calculation is
\[
 A\wedge d_MA
 =dz\wedge\bigl[-\eta\partial\eta
       -2\kappa^{-1}\eta D\Phi-\kappa^{-1}D(\Phi\eta)\bigr].
\]
The total derivative vanishes under \(i^*\eta=0\).
The linear transformation includes the ghosts and antifields and
identifies the linear gauge complexes with abelian Chern--Simons
theory in this boundary polarization. Its inverse parameter is the
reason this particular identification excludes \(\kappa=0\).

\subsection{Contracting every BV component jet}

Return to the polynomial parameter ring \(\mathbb k_\kappa\).
The\phantomsection\label{pcsm:def:h373-holomorphic-cotangent-complex} holomorphic cotangent complex for the same operator is
\begin{equation}\label{pcsm:local}
 B_{\mathrm{hol}}=
 \mathbb k_\kappa[u_0,u_1,\ldots]\otimes\Lambda(c_0,c_1,\ldots),
 \qquad du_n=\kappa c_{n+1},\quad dc_n=0.
\end{equation}
Here \(|u_n|=0\) and \(|c_n|=1\).

To compare it with the full fields, use
\(Q\eta=D\eta\) and \(Q\Phi=D\Phi+\kappa\partial\eta\).
When \(Q\) crosses \(dt\) or \(d\bar z\), it acquires a minus sign.
The eight component equations are therefore
\begin{equation}\label{pcsm:eq:component-Q}
\begin{aligned}
 Qc&=0,& Q\phi&=\kappa\partial c,\\
 Q\eta_t&=-\partial_t c,&
 Q\eta_{\bar z}&=-\partial_{\bar z}c,\\
 Q\eta_{t\bar z}&=\partial_t\eta_{\bar z}
                          -\partial_{\bar z}\eta_t,\\
 Q\phi_t&=-\partial_t\phi-\kappa\partial\eta_t,&
 Q\phi_{\bar z}&=-\partial_{\bar z}\phi
                                   -\kappa\partial\eta_{\bar z},\\
 Q\phi_{t\bar z}&=\partial_t\phi_{\bar z}
              -\partial_{\bar z}\phi_t+\kappa\partial\eta_{t\bar z}.
\end{aligned}
\end{equation}
Extend by commuting with all three derivatives and by the graded
Leibniz rule. This retains the equations and gauge identities in
addition to the transformations of physical fields.

Let\phantomsection\label{pcsm:def:component-jet-module} \(\tau_1,\tau_2,\zeta\) record the undivided jet derivatives
\(\partial_t,\partial_{\bar z},\partial_z\).
The linear jet module is
\[
 V_{\mathrm{jet}}=\mathbb k_\kappa[\tau_1,\tau_2,\zeta]
             \otimes\bigoplus_{I\subseteq\{1,2,3\}}\mathbb C e_I,
 \qquad |e_I|=1-|I|,
\]
with component basis
\[
\begin{array}{c|c@{\qquad}c|c}
 I&e_I&I&e_I\\ \hline
 \varnothing&c&3&\phi\\
 1&-\eta_t&13&-\phi_t\\
 2&-\eta_{\bar z}&23&-\phi_{\bar z}\\
 12&-\eta_{t\bar z}&123&-\phi_{t\bar z}
\end{array}
\]

Index strings are increasing.
Identify this basis with an exterior-index basis in three symbols
\(\theta_1,\theta_2,\theta_3\), with the displayed shifted degrees.
Left contraction \(\iota_j\) removes index \(j\) with its exterior
sign. Then every equation in \eqref{pcsm:eq:component-Q} is contained in
\begin{equation}\label{pcsm:eq:jet-Koszul}
 Q=\tau_1\iota_1+\tau_2\iota_2+\kappa\zeta\iota_3.
\end{equation}

The\phantomsection\label{pcsm:def:polynomial-bv-jets} full polynomial observable algebra is
\(B_{\mathrm{BV,jet}}=\operatorname{Sym}^{\mathrm{gr}}_{
\mathbb k_\kappa}(V_{\mathrm{jet}})\).
Each element uses finitely many jets. There is no largest-jet quotient
and no integration over the point.

Give\phantomsection\label{pcsm:def:jet-projection} a linear generator \(\tau_1^a\tau_2^b\zeta^n e_I\)
transverse weight \(a+b+|I\cap\{1,2\}|\), and add weights on
products. Define \(\pi\) to kill the positive-weight generators and
to send
\(\zeta^ne_\varnothing\mapsto c_n\),
\(\zeta^ne_3\mapsto u_n\).
Its multiplicative section \(\jmath\) sends \(c_n\) to
\(\partial^nc\) and \(u_n\) to \(\partial^n\phi\).

\begin{proposition}[Complete abelian jet reduction]
\label{pcsm:prop:component-jets}
The maps \(\pi:B_{\mathrm{BV,jet}}\to B_{\mathrm{hol}}\) and
\(\jmath:B_{\mathrm{hol}}\to B_{\mathrm{BV,jet}}\) are dg algebra
maps with \(\pi\jmath=1\).
There is a degree-minus-one polynomial homotopy \(H\) such that
\begin{equation}\label{pcsm:eq:jet-contraction}
 QH+HQ=1-\jmath\pi.
\end{equation}
All three maps commute with holomorphic translation and with scalar
specialization of \(\kappa\).
\end{proposition}

\begin{proof}
The differential preserves transverse weight.
At weight zero, \eqref{pcsm:eq:jet-Koszul} reads
\(Qe_3=\kappa\zeta e_\varnothing\), which is precisely
\eqref{pcsm:local}. This proves the assertions about the algebra maps.

On the linear jet module put
\[
 h_0=\theta_1\wedge\partial_{\tau_1}
                      +\theta_2\wedge\partial_{\tau_2}.
\]
The wedges act on the exterior index, not by multiplying observables.
Contraction and wedge give \(Qh_0+h_0Q=N\), where \(N\)
is transverse weight. The term \(\kappa\zeta\iota_3\)
anticommutes with \(h_0\).

Extend \(h_0\) as a graded derivation \(H_0\) of the polynomial
algebra. Its anticommutator with \(Q\) is a derivation, and the
generator identity proves \(QH_0+H_0Q=N\) on every product.
Both operators preserve weight.

Define \(H=N^{-1}H_0\) at
positive weight and \(H=0\) at weight zero.
This gives \eqref{pcsm:eq:jet-contraction}.
Every sum is finite on a polynomial. Only positive integers are
inverted. The formulas are polynomial in \(\kappa\) and commute
with the holomorphic translation derivation.
\end{proof}

For an affine Hamiltonian operator the nonlinear differential admits
the corresponding contraction over \(\mathbb C[\kappa]\)
(Theorem~\ref{bcz:thm:jets}).
For a general Hamiltonian operator, the comparison must still be
constructed with its nonlinear differential. The independently
defined complex \(\mathcal P_H\) by itself supplies no such map.

\subsection{Absolute cohomology and translation}

The transverse reduction is valid at zero level, but the further
contraction of holomorphic jets requires a nonzero level.

\begin{proposition}[Absolute local cohomology]\label{pcsm:localcohom}
After specialization at \(\kappa\in\mathbb C^\times\),
\(H^*(B_{\mathrm{hol}})=\Lambda_{\mathbb C}(c_0)\).
\end{proposition}

\begin{proof}
Let \(N\) count occurrences of \(u_n,c_{n+1}\), leaving \(c_0\)
uncounted. On positive count, use left odd derivatives to define
\[
 h=\frac1{\kappa N}\sum_{n\geq0}u_n\partial_{c_{n+1}}.
\]
Set \(h=0\) on \(1,c_0\).
The anticommutator of
\(\sum_n c_{n+1}\partial_{u_n}\) and
\(\sum_n u_n\partial_{c_{n+1}}\) is the counting operator \(N\).
Thus \(dh+hd=1-ip\), where \(p\) kills all counted variables
and \(i\) includes \(\Lambda(c_0)\).
No infinite sum acts on an individual polynomial.
\end{proof}

The surviving \(c_0\) is the constant ghost in the absolute
Lie-algebraic complex. A theory taking invariants under a specified
global gauge group can use a different complex. In particular,
constant reductive-group invariants need not be represented by this
ghost \cite{pcsm:CDG}*{\S6.2.1}.
The smooth quantum observables and their supported products remain
\eqref{pcsm:quantobs}.

The\phantomsection\label{pcsm:def:constant-operator-complex} same polynomial argument handles any nonzero odd polynomial
\(p(T)\in\mathbb C[T]\) in translation.
For the constant operator \(p(\partial)\), let \(B_p\) have
\(du_n=p(\partial)c_n\), \(dc_n=0\).
Its linear map from even to odd generators is multiplication by
\(p(T)\) on \(\mathbb C[T]\), hence injective.

After splitting off its image, contracting those paired generators leaves
\[
 H^*(B_p)=\Lambda\bigl(\mathbb C[T]/p(T)\bigr).
\]
For \(p(T)=\kappa T\) there is one odd class.
At \(p=0\) the differential vanishes and every generator survives.
This argument concerns polynomial local complexes, without an analytic
quantization claim for higher-order operators.

Translation distinguishes the zero-level fiber too.
The degree-minus-one derivation
\(K=\sum_{n\geq0}u_n\partial_{c_n}\) satisfies
\[
 dK+Kd=\kappa\partial.
\]
Both sides send \(u_n\) to \(\kappa u_{n+1}\) and \(c_n\) to
\(\kappa c_{n+1}\), which proves the identity on the whole algebra.
At nonzero level, \(K/\kappa\) is a null-homotopy of translation.
At zero level, translation takes the nonzero class \(c_0\) to
\(c_1\), so no such homotopy exists.

\section{Interior insertions act through contractions}

The small point-local cohomology does not make the action of supported
bulk cochains trivial.
An eliminated normal variable can contract with a boundary insertion
even when its own projection is zero. A polynomial pair shows the
mechanism before we apply it to the collar maps.

\subsection{One normal pair}

Over\phantomsection\label{pcsm:def:normal-pair} \(R\), take \(\mathcal N=R[x,u]\otimes\Lambda(v)\), with
\(|x|=|u|=0\) and \(|v|=-1\), and define
\[
 D_0=(u-\hbar\kappa\partial_x)\partial_v,
 \qquad
 p^q=\left.\exp(\hbar\kappa\partial_u\partial_x)
                                      \right|_{u=v=0}.
\]
The exponential is finite on polynomials.
For \(i(f)=f(x)\), direct differentiation gives
\(D_0^2=0\), \(p^qD_0=0\) and \(p^qi=1\).

More explicitly,
\[
 p^q(u^ax^b)=
 \begin{cases}
  (\hbar\kappa)^a\dfrac{b!}{(b-a)!}x^{b-a},&a\leq b,\\
  0,&a>b.
 \end{cases}
\]
Thus \(p^q(u)=0\), but \(p^q(ui(f))=\hbar\kappa\partial_xf\).
The exact element \(D_0v=u\) has a nonzero operator image at
nonzero level.

This operation cannot be a cochain map into
\(\operatorname{End}_R(R[x])\) with zero differential.
For a geometric cochain action we must keep the tangential complex
and the actual supported boundary inclusion.

\subsection{The supported bulk action}

Use the collar profile of \eqref{pcsm:normalI} and its quantum maps.
Let \(U\subset V\times(a,b)\), where
\(0<\varepsilon<a<b\).
An observable in \(U\) and the image of \(i^q\) have disjoint
normal supports. Extension by zero, their factorization product and
the quantum projection define
\begin{equation}\label{pcsm:beta}
 \begin{split}
 \beta_{U;V}:\mathcal O^q(U)&\longrightarrow
       \operatorname{Hom}^\bullet_R
                    (\mathcal A_\kappa(V),\mathcal A_\kappa(V)),\\
 \beta_{U;V}(f)(a_0)&=
       p^q\bigl(\operatorname{ext}_U(f)\,i^q(a_0)\bigr).
 \end{split}
\end{equation}
The target consists of continuous cochain operations, with
Hom differential \(T\mapsto dT-(-1)^{|T|}Td\).

The two quantum comparison maps are cochain maps, and disjointness
removes the cross BV terms in the middle product.
Hence \(\beta_{U;V}\) is a cochain map.

Its mixed term can be nonzero when the linear projection vanishes.
Choose\phantomsection\label{pcsm:def:mixed-collar-test} a unit-integral normal form \(\rho_{\mathrm{int}}dt\)
above the boundary profile and compact tangential tests
\(\phi\in\Omega_c^{0,0}(V)\),
\(\beta\in\Omega_c^{0,1}(V)\), with
\(\tau_{\mathbb C}(\phi\beta)=1\).
In the shifted conventions set
\[
 u=\zeta^{-1}\kappa s(\rho_{\mathrm{int}}dt\,\phi e_\Phi),
 \qquad v=s\beta.
\]
Then \(p(u)=0\), whereas on the boundary profile
\(h_Cu=\zeta^{-1}\kappa s(\phi e_\Phi)\).

Since \(h_Ci=0\) and \(ip(u)=0\), the symmetrized tensor
homotopy on \(u\otimes i(v)\) is \(h_Cu\otimes i(v)\).
Using \(i(v)=-\zeta sI(\beta e_\eta)\) and the trace sign,
\[
 P_2(u,i(v))=-b(h_Cu,i(v))=\kappa.
\]
Consequently
\[
 \beta(u)(v)=\hbar\kappa,\qquad \beta(u)(1)=0.
\]

This is the contraction lost by ordinary restriction.
The test \(u\) is not assumed closed. The calculation is an
operation on cochains, not a new degree-zero bulk cohomology class.

For homogeneous linear \(u\) and boundary generators \(v_j\),
the whole operation is
\[
 \begin{split}
 \beta(u)(v_1\cdots v_r)
 &=p(u)v_1\cdots v_r\\
 &\quad+\hbar\sum_{j=1}^r
 (-1)^{|u|\sum_{\ell<j}|v_\ell|}P_2(u,i(v_j))
                         v_1\cdots\widehat v_j\cdots v_r.
 \end{split}
\]
The second term retains its factor \(\hbar\) over every \(R_N\)
and over \(\mathbb k_\kappa[[\hbar]]\).
It cannot be divided by \(\hbar\) on either coefficient ring.

Ordered normal supports also give a composition homotopy.
Suppose \(f\) is supported above \(g\).
The normal primitive of \(g\), together with its boundary profile,
stays below \(f\), so multiplication by \(f\) remains a cochain
map during the contraction.
The joint projected product and successive \(\beta\)-operations
differ by inserting \(1-i^qp^q=dH^q+H^qd\).

Thus their difference is the Hom differential of
\[
 K(f,g,a_0)=(-1)^{|f|}p^q\bigl(
  \operatorname{ext}(f)\,
  H^q(\operatorname{ext}(g)\,i^q(a_0))\bigr).
\]

The sign moves the degree-minus-one homotopy past \(f\).
This supplies the ordered composition comparison with its support
condition. Coherence under arbitrary motion of supports needs the
corresponding parameter-dependent operations.

\section{Gauge fields and the gravitational sector}
\label{pcsm:gravity}

The Gaussian example used a constant response operator.
For currents of a Lie algebra, the field-dependent term expresses
the commutator of gauge transformations.
Let \(\mathfrak g\) be a finite-dimensional complex Lie algebra
with nondegenerate invariant symmetric form \(b\).
Choose a basis \(t_a\), write
\([t_a,t_b]=f_{ab}^{\ d}t_d\), and let \(J_a\) be its currents.
Their affine bracket is
\[
 \{J_a{}_{\lambda}J_b\}=f_{ab}^{\ d}J_d
                              +\kappa b(t_a,t_b)\lambda.
\]
The convention \(\{J_a{}_{\lambda}J_b\}=H^{ba}(\lambda)\)
fixes the transpose of the field-linear coefficient.

The Chern--Simons interpretation can be seen in the physical action.
Identify the currents with \(\Phi\in\mathfrak g\) by
\(J_a=b(t_a,\Phi)\), and write \(a=\eta_a t_a\) for the
physical partial connection. The action becomes
\[
 S^{\mathrm{cl}}_H=
 \tau_M\left(b(a,D\Phi)-\tfrac12b(\Phi,[a,a])
                              +\tfrac\kappa2b(a,\partial a)\right).
\]
The Lie bracket includes exterior multiplication of form coefficients.

For \(\kappa\ne0\), set
\(\mathscr A=-a+\kappa^{-1}\Phi\,dz\) and normalize
\[
 \operatorname{CS}(\mathscr A)=\frac1{2\pi i}\int_M
 \left(b(\mathscr A,d_M\mathscr A)
                +\tfrac13b(\mathscr A,[\mathscr A,\mathscr A])\right).
\]
The quadratic calculation of \eqref{pcsm:csaction} and the three
placements of \(\Phi\,dz\) in the cubic term give
\[
 -\frac\kappa2\operatorname{CS}(\mathscr A)
       =S^{\mathrm{cl}}_H+\tfrac12\tau_M D\,b(\Phi,a).
\]
The boundary term vanishes for \(i^*a=0\).
At \(\kappa=0\), the displayed first-order action instead has
the holomorphic--topological BF form.
This is the local classical affine realization
\cite{pcsm:KZ}*{\S3.1}.

A comparison of full BV theories also retains the symplectic form,
ghosts and antifields. With a full-superfield map \(F\), coefficient
\(c_{\mathrm{CS}}\) and local two-form \(B\), its normalization
is expressed by
\[
 F^*\omega_{\mathrm{CS}}=\omega_{\mathrm{PV}},\qquad
 F^*(c_{\mathrm{CS}}\operatorname{CS})=S_H+\tau_M DB.
\]
A numerical identification with a gravitational coupling or a quantum
boundary level must use these conventions and the corresponding
quantum comparison.

\subsection{The Virasoro interaction}

A stress field supplies a nonlinear scalar example.
For an even generator \(T\) and \(c\in\mathbb C\), take
\[
 \{T_\lambda T\}=(\partial+2\lambda)T+\frac c{12}\lambda^3,
 \qquad
 H_T=T'+2T\partial+\frac c{12}\partial^3.
\]
Here \(T'=\partial T\) multiplies the input, while the other
displayed derivatives act on it.
The term \(\tfrac12\eta H_T\eta\) gives the Virasoro sigma-model
interaction directly from the classical bracket.
The Virasoro and \(\mathcal W\) examples of
\cite{pcsm:KZ}*{\S\S3.2--3.5} relate this construction to
three-dimensional gravitational gauge data.

The role of a conformal element can already be stated classically.
Its zero mode acts as holomorphic translation.
The classical translation is a gauge symmetry under the
coordinate-action hypotheses of \cite{pcsm:KZ}*{\S\S2.3 and 4.3}.
These require a diagonal integral or half-integral first-mode grading,
locally nilpotent positive modes and the corresponding spin data. The local translation homotopies are visible in the
polynomial examples below.

Turning this symmetry into a supported quantum topological theory
requires compatible Ward homotopies, radius maps and their coherences.
Neither a global three-manifold theory nor its full topological
field-theory axioms follow from the local symmetry argument.

Write\phantomsection\label{pcsm:def:virasoro-jet-complex} \(u_n=\partial^nT\) and \(c_n=\partial^nc_0\) for the
undivided stress and ghost jets, of degrees zero and one.
On
\[
 P=\mathbb C[u_0,u_1,\ldots]\otimes\Lambda(c_0,c_1,\ldots)
\]
the cotangent differential is the derivation determined by
\[
 du_0=u_1c_0+2u_0c_1+\frac c{12}c_3,
 \qquad dc_0=c_0c_1,
 \qquad [d,\partial]=0.
\]
The nonlinear terms preserve a weight that makes the next calculation
finite on each polynomial.

\begin{proposition}[Polynomial Virasoro cohomology]\label{pcsm:vircohom}
For \(c\ne0\), one has \(H^*(P)=\Lambda(\xi)\), where
\(\xi=[c_0c_1c_2]\) has degree three.
\end{proposition}

\begin{proof}
Give the generators weights \(w(u_n)=n+2\) and \(w(c_n)=n-1\).
The differential preserves weight. Its central part is
\[
 \delta u_n=\frac c{12}c_{n+3},\qquad \delta c_n=0.
\]

At a fixed weight, filter decreasingly by
\(\deg_{\mathrm{coh}}+\#u\), where \(\#u\) counts even
factors. The operator \(\delta\) preserves this number and
\(d-\delta\) increases it by one.
This is a finite filtration: only \(c_0\) has negative weight,
only \(c_1\) has weight zero, and both occur at most once.
All other variables have positive weight.

The pairs \((u_n,c_{n+3})\) contract for \(\delta\), leaving
\(\Lambda(c_0,c_1,c_2)\).
The next differential is
\[
 dc_0=c_0c_1,\qquad dc_1=c_0c_2,\qquad dc_2=c_1c_2.
\]
It is an isomorphism from degree one to degree two and vanishes
on degree two. Only the unit and the degree-three product survive.
No later degree-one differential can join their degrees.
The original product \(c_0c_1c_2\) is closed, so it represents
the surviving class in the convergent weightwise calculation.
\end{proof}

The explicit polynomial comparison in
Theorem~\ref{vir:contraction} explains what the spectral sequence
has eliminated. The three residual ghosts form the
\(\mathfrak{sl}_2\) gauge complex. Its section corrects the third
ghost by a nonlinear stress term, and its finite homotopy contracts
the other pairs. Those same corrected coordinates construct the
oper connection used for finite-dimensional lines below.

\subsection{Integrating a holomorphic density changes the complex}

Let\phantomsection\label{pcsm:def:integrated-density-complex} \(F=P/\partial P\) be the complex of translation-invariant
holomorphic differential-polynomial densities modulo total derivatives.
It has neither the support nor the carrier of all three-dimensional
local functionals.

The kernel of \(\partial:P\to P\) is \(\mathbb C\).
In a nonconstant polynomial, the next highest-jet coefficient
forces independence of the highest jet.
This applies to even and exterior variables alike. Descending through the jet orders proves the claim.
Thus there is an exact sequence of complexes
\[
 0\longrightarrow P/\mathbb C\xrightarrow{\partial}P
                 \longrightarrow F\longrightarrow0.
\]

Left\phantomsection\label{pcsm:def:virasoro-translation-homotopy} exterior differentiation \(K=\partial/\partial c_0\) has
degree \(-1\). The generator formulas give
\[
 (dK+Kd)u_n=u_{n+1},\qquad
 (dK+Kd)c_n=c_{n+1}.
\]
Both sides are derivations, so \(dK+Kd=\partial\) on \(P\).
Since \(K1=0\), this is also a homotopy of the displayed map
\(P/\mathbb C\to P\).

Its cone therefore splits, as a complex, and
\[
 F\simeq P\oplus(P/\mathbb C)[1].
\]
For \(c\ne0\), its cohomology has one dimension in each of
degrees \(0,2,3\), with density representatives
\(1,c_1c_2,c_0c_1c_2\).
The degree-two representative is closed in the quotient because
\(d(c_1c_2)=\partial(c_0c_1c_2)\).

Multiplication and the entire Schouten \(\lambda\)-bracket do not
descend to this quotient. Write \(\int f\) for the class of \(f\).
Its integrated zero-mode bracket does descend:
\[
 \{\!\int f,\int g\!\}
       =\int\{f_\lambda g\}_{\mathrm{Sch}}\big|_{\lambda=0}.
\]
The two sesquilinearity identities make any change by a total
derivative vanish in this expression.
The integrated complex and the full polynomial complex consequently
retain different operations as well as different cohomology groups.

\section{Line operators and the full local symmetry}
\label{pcsm:sec:witt-lines}

A line coupled to every holomorphic jet of the Virasoro gauge field
must respond to every derivative of its gauge parameter.
This requires an action of all local vector fields on its coefficient
module. The three-dimensional subalgebra \(\mathfrak{sl}_2\)
does not suffice for a field-independent coupling.
We shall construct an infinite polynomial module that carries the
full action, then a different finite-dimensional transport that uses
the stress field itself.

Let\phantomsection\label{pcsm:def:witt-algebra}
\[
 \mathfrak w_+=\bigoplus_{n\geq-1}\mathbb C\ell_n,
 \qquad \ell_n=-z^{n+1}\partial_z,
 \qquad [\ell_n,\ell_m]=(n-m)\ell_{n+m}.
\]
The bracket of two \(\ell_{-1}\)'s is zero.
Its subalgebra spanned by \(\ell_{-1},\ell_0,\ell_1\) is
\(\mathfrak{sl}_2\), with
\(E=\ell_1\), \(F=-\ell_{-1}\), \(H=-2\ell_0\).
In the two-dimensional representation, \(\ell_0\) has weights
\(-1/2,1/2\). The relation
\([\ell_0,\ell_2]=-2\ell_2\) forces \(\ell_2=0\), since
no pair of weights differs by two.
But \([\ell_2,\ell_{-1}]=3\ell_1\) would then force the
nonzero raising operator to vanish.

\begin{proposition}\label{pcsm:prop:witt-finite}
Every finite-dimensional representation of \(\mathfrak w_+\) is
trivial. A nontrivial finite-dimensional \(\mathfrak{sl}_2\)
representation therefore has no extension to \(\mathfrak w_+\)
on its original carrier.
\end{proposition}

\begin{proof}
The Lie algebra of polynomial vector fields is simple.
Indeed, a nonzero ideal contains \(f(z)\partial_z\) for a
nonzero polynomial \(f\).
Repeated brackets with \(\partial_z\) give a nonzero constant
multiple of \(\partial_z\).
Bracketing it with \(z^{m+1}\partial_z\) gives
\((m+1)z^m\partial_z\), for every \(m\geq0\), so the ideal
contains all polynomial vector fields.
A map from this infinite-dimensional Lie algebra to the
finite-dimensional vector space \(\operatorname{End}_{\mathbb C}(W)\)
has a nonzero kernel. Simplicity makes that kernel the whole algebra.
\end{proof}

The polynomial Neveu--Schwarz superalgebra has even part
\(\mathfrak w_+\) and odd generators
\(g_r\), \(r=-1/2,1/2,3/2,\ldots\), with graded brackets
\[
 [\ell_n,g_r]=(n/2-r)g_{n+r},\qquad
 [g_r,g_s]=2\ell_{r+s}.
\]
A zero-coefficient term with an index outside this range is omitted.
Every finite-dimensional superrepresentation is trivial as well:
the even action vanishes by the proposition, and
\([\ell_0,g_r]=-r g_r\) then kills every odd generator.

Algebraic induction enlarges a spin space, but need not make an
infinite jet act on it. A \(\mathfrak w_+\)-module \(M\) is
\emph{restricted} if, for every \(v\in M\), one has
\(\ell_n v=0\) for all sufficiently large \(n\), with the bound
allowed to depend on \(v\). This condition makes a formal sum of
positive modes act on each state by a finite sum.

\begin{proposition}\label{pcsm:prop:induced-restricted}
For \(r\geq1\), put \(V_r=\operatorname{Sym}^r\mathbb C^2\) and
\[
 I_r=U(\mathfrak w_+)\otimes_{U(\mathfrak{sl}_2)}V_r.
\]
The map \(v\mapsto1\otimes v\) is injective, and for every
algebraic \(\mathfrak w_+\)-module \(M\) there is a natural
isomorphism
\[
 \operatorname{Hom}_{\mathfrak w_+}(I_r,M)
 \simeq\operatorname{Hom}_{\mathfrak{sl}_2}
                 (V_r,M|_{\mathfrak{sl}_2}).
\]
The module \(I_r\) is not restricted. If \(M\) is restricted,
both displayed spaces are zero.
\end{proposition}

\begin{proof}
Order a basis of the span of \(\ell_n\), \(n\geq2\), before
the three \(\mathfrak{sl}_2\) basis vectors. The
Poincar\'e--Birkhoff--Witt theorem makes \(U(\mathfrak w_+)\)
free as a right \(U(\mathfrak{sl}_2)\)-module, with the ordered
monomials in the first group as a basis. This proves injectivity.
Restriction of a map to \(1\otimes V_r\) has inverse
\(f\mapsto(u\otimes v\mapsto u f(v))\), proving the adjunction.
For every \(N\geq2\) and \(v\ne0\), the vector
\(\ell_N(1\otimes v)=\ell_N\otimes v\) is a nonzero PBW vector.
Thus \(I_r\) is not restricted.

Let \(W\) be the image of an \(\mathfrak{sl}_2\)-linear map
\(V_r\to M\), with \(M\) restricted. A finite basis of \(W\)
has a common bound, so \(\ell_NW=0\) for some \(N\geq2\).
Since \(\ell_{-1}W\subset W\), the relation
\([\ell_N,\ell_{-1}]=(N+1)\ell_{N-1}\) implies
\(\ell_{N-1}W=0\). Descending gives \(\ell_1W=0\).
The remaining two brackets then kill \(\ell_0W\) and
\(\ell_{-1}W\). The irreducible nontrivial module \(V_r\)
has no nonzero image with trivial \(\mathfrak{sl}_2\) action,
so \(W=0\).
\end{proof}

An infinite highest-weight string provides a restricted replacement.

\subsection{A polynomial module for every local vector field}

Fix\phantomsection\label{pcsm:def:polynomial-witt-module} \(a\in\mathbb C\) and take \(M_a=\mathbb C[x]\) in
degree zero. With \(d_x=\partial_x\) the ordinary, undivided
derivative, define
\begin{equation}\label{pcsm:eq:witt-polynomial}
 L_{-1}=x,\qquad
 L_n=xd_x^{n+1}+(n+1)a\,d_x^n\quad(n\geq0).
\end{equation}
The derivative orders make sufficiently high modes vanish on each
fixed polynomial, while \(L_{-1}\) raises its degree by one.

\begin{proposition}\label{pcsm:prop:witt-polynomial}
The assignments \(\ell_n\mapsto L_n\) give a faithful
\(\mathfrak w_+\)-representation on \(M_a\).
Its \(\mathfrak{sl}_2\)-restriction is the Verma module with
highest vector \(1\) and highest \(H\)-weight \(-2a\).
The action extends continuously to
\(\widehat{\mathfrak w}_+=\prod_{n\geq-1}\mathbb C\ell_n\),
with the product topology of discrete coefficient lines. The target
has pointwise evaluation on discrete polynomials.
\end{proposition}

\begin{proof}
Use \(d_x^r x=xd_x^r+r d_x^{r-1}\) to obtain, for \(n,m\geq0\),
\[
 [L_n,L_m]=(n-m)
       \bigl(xd_x^{n+m+1}+(n+m+1)a\,d_x^{n+m}\bigr),
 \qquad [L_n,x]=(n+1)L_{n-1}.
\]
These are all the required relations.
Since \(L_{-1}\ne0\), simplicity implies faithfulness.

The vectors \(F^k1=(-x)^k\) form a basis, with
\(E1=0\) and \(H1=-2a\). Hence the restricted module has
exactly the Verma basis and action.
Explicitly, for \(0\leq n\leq k\),
\[
 L_nx^k=(k)_n\bigl(k-n+(n+1)a\bigr)x^{k-n},
 \qquad (k)_n=k(k-1)\cdots(k-n+1),\quad(k)_0=1.
\]
For \(n>k\) the value is zero.

A formal sum of modes thus has a finite value on every state.
In a composite, each factor raises degree by at most one, so the
relevant double sum remains finite.
The coefficient of \(\ell_j\) in the formal Lie bracket also
has finitely many terms, because both indices are at least \(-1\).
The polynomial bracket identity extends coefficientwise, and the
same bounds prove continuity.
\end{proof}

At \(a=-m/2\), \(m\geq0\), the irreducible
\(\mathfrak{sl}_2\) quotient is \(M_a/(x^{m+1})\).
For \(m\geq1\), its defining submodule fails to be invariant
under the full vector-field algebra:
\begin{equation}\label{pcsm:eq:witt-quotient-defect}
 L_2x^{m+1}=-\frac{m(m+1)(m+2)}2\,x^{m-1}.
\end{equation}
The trivial quotient at \(m=0\) does extend, with every mode zero.
Thus the polynomial module supplies the line coupling, whereas its
nontrivial finite-dimensional quotient loses a required relation.

\subsection{Formal transport and endpoint covariance}

Let\phantomsection\label{pcsm:def:virasoro-gauge-field} \(\mu=\mu_tdt+\mu_{\bar z}d\bar z\) be the physical
Virasoro gauge field and fix a vertical line at \(z=z_0\).
For an even parameter \(\epsilon\), its time component transforms as
\begin{equation}\label{pcsm:eq:witt-gauge}
 \delta\mu_t=-\partial_t\epsilon
                   +\epsilon\partial_z\mu_t-\mu_t\partial_z\epsilon.
\end{equation}
The local jet coupling of \cite{pcsm:KZ}*{\S3.2.2} is now defined
on the actual module \(M_a\):
\[
 A_a(t)=\sum_{n\geq0}
     \frac{\partial_z^n\mu_t(t,z_0,\bar z_0)}{n!}L_{n-1},
 \qquad
 B_a(t)=\sum_{n\geq0}
     \frac{\partial_z^n\epsilon(t,z_0,\bar z_0)}{n!}L_{n-1}.
\]
Each sum is finite on a polynomial.
Taylor coefficients identify the nonlinear term of
\eqref{pcsm:eq:witt-gauge} with
\([\widehat\mu_t,\widehat\epsilon]\) in the formal vector-field
algebra, giving
\begin{equation}\label{pcsm:eq:witt-connection-gauge}
 \delta A_a=-\partial_tB_a+[A_a,B_a].
\end{equation}

For formal time expansions \(A_a(t)=\sum_{j\geq0}A_jt^j\),
define \(U_a(t)=\sum_{k\geq0}U_kt^k\) by
\begin{equation}\label{pcsm:eq:witt-transport}
 U_0=1,\qquad
 U_{k+1}=\frac1{k+1}\sum_{j=0}^k A_jU_{k-j}.
\end{equation}
This is the unique solution of \(\partial_tU_a=A_aU_a\) with
identity initial value. It is invertible as a formal operator series.
On each state its value lies in \(M_a[[t]]\): every time coefficient
uses finitely many polynomial operations, with no completion in \(x\).

\begin{proposition}\label{pcsm:prop:witt-transport}
The transport is gauge covariant:
\[
 \delta U_a(t)=-B_a(t)U_a(t)+U_a(t)B_a(0).
\]
It transports endpoint states transforming by \(-B_a\).
\end{proposition}

\begin{proof}
Differentiate the right side, substitute
\eqref{pcsm:eq:witt-connection-gauge}, and use
\(\partial_tU_a=A_aU_a\).
It satisfies the linearized transport equation with zero initial
variation. The coefficient recursion gives uniqueness.
\end{proof}

This gives a formal open line on an infinite coefficient module.
An analytic evolution on that module, or a trace closing the line,
requires its own domain. Finite-dimensional transport instead arises
from a connection depending on the stress field.

\subsection{Finite-dimensional transport from the oper connection}

For\phantomsection\label{pcsm:def:local-oper} \(c\ne0\), put \(q=6T/c\), with primes denoting
\(\partial_z\). The equation \(\psi''+q\psi=0\) can be
written on the two-component column \((\psi,\psi')^{\mathsf t}\).
Its geometric form is an oper.
Over a commutative differential algebra \((\mathcal O,\partial_z)\),
let \(\mathcal E\) be a free rank-two module with trivialized
determinant. A local \(\mathrm{SL}_2\) oper consists of a
determinant-preserving \(z\)-connection \(\nabla\) and a rank-one
direct summand \(\mathcal L\subset\mathcal E\), for which
\[
 \mathcal L\xrightarrow{\nabla}
         (\mathcal E/\mathcal L)\otimes\Omega_z^1,
 \qquad \Omega_z^1=\mathcal O\,dz,
\]
is an isomorphism. Projection to the quotient makes this map
\(\mathcal O\)-linear.

Take\phantomsection\label{pcsm:def:oper-line-basis} \(\mathcal E=\mathcal O^{\oplus2}\),
\(\mathcal L=\mathcal O e_2\), and
\(\nabla=(\partial_z-A_z)dz\), with
\begin{equation}\label{pcsm:eq:oper-line-matrices}
 A_z=\begin{pmatrix}0&1\\-q&0\end{pmatrix},\qquad
 E_q(f)=\begin{pmatrix}f'/2&-f\\qf+f''/2&-f'/2\end{pmatrix},
 \qquad A_t=E_q(\mu_t).
\end{equation}
The induced map sends \(e_2\) to \(-e_1dz\), so the oper
condition holds. These are the matrices of
Chapter~\ref{rc:virasoro-oper} with its variables \(U=-q\) and
\(\nu=-\mu\).
For time-dependent coefficients, let \(\partial_t\) commute with
\(\partial_z\). The Virasoro transformation and the time equation
of motion become
\[
 \delta q=\epsilon q'+2\epsilon'q+\tfrac12\epsilon''',
 \qquad
 \mathcal F=\partial_tq+\mu_tq'+2\mu_t'q
                                      +\tfrac12\mu_t'''=0.
\]

\begin{proposition}\label{pcsm:prop:oper-line}
On a classical Virasoro sigma-model solution,
\eqref{pcsm:eq:oper-line-matrices} gives gauge-covariant line
transport in every finite-dimensional \(\mathfrak{sl}_2\)
representation. With \(L=E_q(\epsilon)\), the exact off-shell
identities are
\begin{align}
 \partial_tA_z-\partial_zA_t+[A_z,A_t]
 &=\begin{pmatrix}0&0\\-\mathcal F&0\end{pmatrix},
                              \label{pcsm:eq:oper-line-curvature}\\
 \delta A_z&=-\partial_zL+[A_z,L],\notag\\
 \delta A_t&=-\partial_tL+[A_t,L]
                 +\begin{pmatrix}0&0\\\epsilon\mathcal F&0\end{pmatrix}.
                              \label{pcsm:eq:oper-line-gauge}
\end{align}
\end{proposition}

\begin{proof}
Substitute the displayed transformation of \(q\) and
\eqref{pcsm:eq:witt-gauge} in the matrices.
Their diagonal and upper off-diagonal defects cancel.
The lower entries are respectively \(-\mathcal F\), zero and
\(\epsilon\mathcal F\), proving the identities before imposing
the equations of motion.

On a solution, applying any representation
\(\rho:\mathfrak{sl}_2\to\operatorname{End}_{\mathbb C}(W)\)
therefore gives
\[
 \delta\rho(A_t)=-\partial_t\rho(L)+[\rho(A_t),\rho(L)].
\]
Its transport has the endpoint covariance of
Proposition~\ref{pcsm:prop:witt-transport}.
For smooth coefficients on a compact time interval it is the
ordinary finite matrix differential equation.
In the fundamental representation \(\operatorname{Tr}A_t=0\),
so transport preserves the determinant.
\end{proof}

The dependence of \(E_q(\epsilon)\) on \(q\) is essential:
its time covariance uses the stress equation.
This finite-dimensional line and the polynomial jet line have
different coefficient actions. Both give specified endpoints on
which bulk insertions can act. The finite one does not require an
extension of a fixed \(\mathfrak{sl}_2\)-module to
\(\mathfrak w_+\).

\subsection{From flat connections to an Einstein metric}

To\phantomsection\label{pcsm:def:real-cartan-data} pass from gauge fields to a metric, impose a real form and
nondegeneracy. In Lorentzian signature, take
\(\mathfrak{sl}_2(\mathbb R)\)-valued one-forms \(\omega,e\)
and a length \(\ell>0\).
For a connection \(A\), write
\(F(A)=d_MA+\tfrac12[A,A]\), and put
\(d_\omega e=d_Me+[\omega,e]\).
The two connections \(A^\pm=\omega\pm e/\ell\) are flat exactly
when
\begin{equation}\label{pcsm:flat-gravity}
 F(A^\pm)=F(\omega)+\frac1{2\ell^2}[e,e]
                                  \pm\ell^{-1}d_\omega e=0.
\end{equation}

Adding and subtracting gives the constant-curvature and torsion-free
Cartan equations. If \(e\) is invertible, \(\omega\) is the compatible spin connection.
The Einstein metric
\[
 g=2\operatorname{Tr}(e\otimes e)
\]
has Lorentzian signature and sectional curvature \(-\ell^{-2}\).
The trace is in the defining representation.

This is the classical Chern--Simons description
\cites{hol:AT,hol:WittenGravity}.
Flat connections with degenerate \(e\) do not define a metric
in this chart. Euclidean fillings require their corresponding
reality conditions in place of two independent real connections.

Chapter~\ref{filling:chapter} constructs Einstein collars and
complete fillings with boundary data.
The integral affine--ghost reduction of Chapter~\ref{rc:brst}
maps to physical currents under the renormalized current and Ward
assumptions of Theorem~\ref{brst-thm:physical}.
Connection elimination also has determinant and contact terms,
and boundary-domain comparisons have analytic hypotheses
(Chapters~\ref{rc:ec}--\ref{rc:boundary}).
A quantum gravitational identification must preserve these data
along with the classical stress bracket.

\section{Two constructions to compare}
\label{pcsm:reconstruction}

The Poisson vertex bracket determines a classical bulk action. The
Gaussian model also has a quantum boundary complex and a supported
bulk action. Reconstruction asks whether boundary conditions and
the operators between them determine that bulk.
A bulk insertion acts on those objects and morphisms compatibly
with their compositions and with other bulk insertions.
For a unital dg algebra \(A\), the derived algebraic center is
\[
 C^\bullet(A,A)=R\operatorname{Hom}_{A\otimes A^{\mathrm{op}}}(A,A),
\]
with its Hochschild operations (Chapter~\ref{rc:holography}).
The holomorphic question additionally retains translation, singular
collisions and support.

These requirements already separate two kinds of motion in the free
example. On its boundary complex, holomorphic translations cannot
be the unary operations of an ordinary topological moving-disks
action (Proposition~\ref{ca:prop:translation}).
The proposed universal object is instead the diagonal center of
Conjecture~\ref{rv:intro:conj:ht-centre}, with its unital supported
mixed factorization structure. Existence and terminality in that
specified action category remain to be proved.
The map \eqref{pcsm:beta} supplies an explicit cochain action and
its ordered composition homotopy on which to test this comparison.

Choosing one boundary object introduces a further generation question.
Let \(P\) be a right dg \(A\)-module, replaced by a homotopically
projective representative, and put \(B=\operatorname{End}_A(P)\).
It models \(R\operatorname{Hom}_A(P,P)\), and \(P\) is a
\(B\)--\(A\) bimodule. The interface gives the adjunction on
right derived module categories
\[
 -\otimes_B^{\mathbb L}P:\operatorname{Mod}_B
 \rightleftarrows\operatorname{Mod}_A:
 R\operatorname{Hom}_A(P,-).
\]
If \(P\) is compact and detects acyclic modules, it is a compact
generator and the adjunction is an equivalence
(Proposition~\ref{hol:generator}).
Transport of a holomorphic bulk action must further intertwine
the supported collision operations and translations.

An augmentation object can detect only a completion.
For \(A=\mathbb C[x]\) and \(P=\mathbb C\) with \(x\) acting
by zero, the first endomorphism algebra is modeled by
\(\Lambda(\eta)\), \(|\eta|=1\).
Its ungraded derived double centralizer is \(\mathbb C[[x]]\),
and the natural map is
\(\mathbb C[x]\to\mathbb C[[x]]\)
(Proposition~\ref{rbt:prop:double-dual}).
The choice of boundary object has retained a formal neighborhood.
Recovering the full theory requires the corresponding generation
and operation comparisons.

Holomorphic bulk--boundary structures occur in the twisted holography
of Costello and Gaiotto~\cite{hol:CG}.
Costello--Paquette formulate a universal chiral defect coupling
through Koszul duality, with additional hypotheses on the free limit
\cite{hol:CP}*{\S6.3}. Categorical boundary constructions also
produce proposed bulk descriptions in \cite{hol:categorical}.

For a prescribed Poisson vertex bracket, the sigma model defines
a concrete reconstruction problem.
First, construct its quantization and boundary category.
Then compare the supported action with the proposed center and
determine which bulk sector the boundary probes generate.
An Einstein realization must also carry the independently specified
states and the sewing maps of their fillings.
